\documentclass[11pt]{article}
\usepackage[margin=1in]{geometry}
\usepackage{amsmath,amssymb,amsthm,mathtools}
\usepackage{microtype}
\usepackage[hidelinks]{hyperref}
\usepackage{enumitem}
\usepackage{booktabs}

\newtheorem{theorem}{Theorem}
\newtheorem{proposition}[theorem]{Proposition}
\newtheorem{lemma}[theorem]{Lemma}
\newtheorem{corollary}[theorem]{Corollary}
\theoremstyle{definition}
\newtheorem{definition}[theorem]{Definition}
\newtheorem{remark}[theorem]{Remark}

\newcommand{\A}{\mathbb A}
\newcommand{\C}{\mathbb C}
\newcommand{\Sing}{\operatorname{Sing}}
\newcommand{\Crit}{\operatorname{Crit}}
\newcommand{\CV}{\operatorname{CV}}
\newcommand{\rank}{\operatorname{rank}}
\newcommand{\mult}{\operatorname{mult}}

\title{A Non-Suspension Polynomial Solution in Gao's Five-Dimensional Middle Branch}
\author{Casey Atwell}
\date{18 August 2026}

\begin{document}
\maketitle

\begin{abstract}
For the mixed direction field
\[
\Delta=(1,w_1,w_1^2,w_2)^{\mathsf T},
\]
Gao's five-dimensional middle branch is governed by the polynomial identity
$4cD_0=D_1^2$ and was conjectured to be rigid up to suspension of a four-dimensional sweep. We give explicit polynomial data satisfying the middle-branch equation and prove that the associated unpadded sweep is not polynomially equivalent to any triangular one-coordinate suspension of a four-dimensional Gao sweep. No determinant normalization is imposed on the suspension: polynomial equivalence to our sweep would force its Jacobian determinant to be a nonzero constant multiple of the square of the Gao sweep parameter, after which Gao's general sweep formula and unique factorization force the two possible base-sweep determinant orders. The proof uses the critical-value closure. For our sweep it is the irreducible threefold
\[
(z_1z_3-z_2^2)^3z_4^2-\lambda z_1^6=0,
\qquad
\lambda=\frac{41-38i}{576}\neq0,
\]
which is not ambient-polynomially equivalent to a cylinder. We formulate the theorem independently of the undefined phrase ``equivalent to a suspension'' in Gao's Open Problem 4.8. Under the source-target polynomial-equivalence interpretation and the triangular one-coordinate suspension class defined here, Open Problem 4.8 therefore has a negative answer.
\end{abstract}

\section{Context and statement of the result}

Gao studies polynomial tangent sweeps associated with the mixed direction field
\[
\Delta=(1,w_1,w_1^2,w_2)^{\mathsf T}
\]
in dimension five. With potentials $\Gamma=(G_2,G_3,G_4)$ and
\[
H:=G_3-2w_1G_2,
\]
Gao writes the determinant pencil as
\[
\det\bigl(J(\Gamma)+\mu\widetilde P\bigr)
   =D_0+\mu D_1+\mu^2D_2,
\]
where
\begin{align*}
D_2&=-\partial_3H,\\
D_1&=(\partial_2H)(\partial_3G_4)
     -(\partial_3H)(\partial_2G_4)
     -\bigl((\partial_1G_2)(\partial_3G_3)
     -(\partial_1G_3)(\partial_3G_2)\bigr),\\
D_0&=\det J(G_2,G_3,G_4).
\end{align*}
The middle branch is normalized by
\begin{equation}
D_2=c\in\C^\times,
\qquad
X_1=-\frac{D_1}{2c},
\qquad
4cD_0=D_1^2.
\label{eq:master}
\end{equation}
This is Gao's equation (12) in arXiv:2608.00222v1 \cite{Gao2026}. Gao's Theorem 4.6 proves that on the characteristic-invariance locus $\mathcal W G_2=0$ every solution collapses to a triangular suspension of a four-dimensional bottom-branch sweep, while Remark 4.7 records another explicit graph-coordinate collapse. Open Problem 4.8 asks whether every polynomial solution of \eqref{eq:master} is equivalent to a suspension of a four-dimensional sweep \cite{Gao2026}.

Throughout, ``five-dimensional middle branch'' follows Gao's padded construction $F_0:\A^5\to\A^5$. The critical-value geometry below concerns its unpadded sweep $S:\A^4\to\A^4$; padding contributes one extra coordinate and one extra factor of $\gamma$ to the Jacobian determinant. Because the phrase ``equivalent to a suspension'' is not formally defined in the statement of Open Problem 4.8, we first state a precise theorem that does not depend on that terminology.

\begin{definition}[Source-target polynomial equivalence]
Two polynomial maps $S_1,S_2:\A^4\to\A^4$ are \emph{polynomially equivalent} if
\[
S_2=\Psi\circ S_1\circ\Phi
\]
for polynomial automorphisms $\Phi,\Psi\in\operatorname{Aut}_{\rm poly}(\A^4)$.
\end{definition}

\begin{definition}[Triangular one-coordinate suspension]
Let $u=(\gamma_T,a,b)$ and let $v$ be a fourth source coordinate. A \emph{triangular one-coordinate suspension of a four-dimensional Gao sweep} is a polynomial map
\[
\Sigma(u,v)=\bigl(T(u),Q(u,v)\bigr):\A^4\to\A^4,
\]
where $T:\A^3\to\A^3$ is the unpadded sweep associated with an arbitrary four-dimensional Gao tangent-sweep datum and $T$ is independent of $v$. No Jacobian-determinant normalization and no determinant branch for $T$ are imposed in the definition.
\end{definition}

\begin{theorem}[Main theorem]\label{thm:main}
There exists an exact polynomial solution of Gao's five-dimensional middle-branch equation \eqref{eq:master} whose associated unpadded sweep $S:\A^4\to\A^4$ is not polynomially equivalent to any triangular one-coordinate suspension of a four-dimensional Gao sweep.
\end{theorem}

\begin{corollary}[Relation to Gao's Open Problem 4.8]\label{cor:gao}
If ``equivalent to a suspension of a four-dimensional sweep'' in Gao's Open Problem 4.8 is interpreted as source-target polynomial equivalence to the triangular one-coordinate suspension class above, then Open Problem 4.8 has a negative answer.
\end{corollary}

\begin{remark}[Scope of the theorem]\label{rem:scope}
Theorem \ref{thm:main} excludes exactly the class defined above: triangular one-coordinate suspensions whose base is an arbitrary four-dimensional Gao sweep, under source-target polynomial equivalence.  It does not claim to classify every object that could informally be called a ``suspension.''  Corollary \ref{cor:gao} is therefore intentionally conditional on interpreting Gao's terminology by this precise class.
\end{remark}

\section{The exact middle-branch witness}

Normalize $c=1$ and $\varphi=0$, and write
\[
x=w_1,\qquad y=w_2,\qquad t=w_3.
\]
Set
\begin{equation}
\boxed{
G_2=yt^2,
\qquad
G_3=2xyt^2-t,
\qquad
G_4=\frac{4i}{3}y^3t^3.
}
\label{eq:witness}
\end{equation}
Then
\[
H=G_3-2xG_2=-t.
\]
This witness lies outside Gao's characteristic-invariance collapse locus. Indeed, for $c=1$ and $\varphi=0$ one has $\mathcal W=\partial_y$, while
\[
\mathcal W G_2=\partial_y(yt^2)=t^2\neq0.
\]
Thus Theorem 4.6 does not apply to the witness.

\begin{proposition}\label{prop:master}
The data \eqref{eq:witness} satisfy Gao's middle-branch conditions with $c=1$:
\[
D_2=1,
\qquad
D_1=4(1+i)y^2t^3,
\qquad
D_0=8iy^4t^6,
\]
and hence $4D_0=D_1^2$ identically.
\end{proposition}

\begin{proof}
From $H=-t$ one has $D_2=-\partial_tH=1$. Since
\[
\partial_yH=0,
\qquad
\partial_tH=-1,
\]
Gao's displayed formula for $D_1$ reduces to
\[
D_1=\partial_yG_4
-\bigl((\partial_xG_2)(\partial_tG_3)
-(\partial_xG_3)(\partial_tG_2)\bigr).
\]
Now $\partial_xG_2=0$, $\partial_xG_3=2yt^2$, $\partial_tG_2=2yt$, and
\[
\partial_yG_4=4iy^2t^3.
\]
Therefore
\[
D_1=4iy^2t^3+4y^2t^3=4(1+i)y^2t^3.
\]
A direct $3\times3$ determinant gives
\[
D_0=\det J(G_2,G_3,G_4)=8iy^4t^6.
\]
Consequently
\[
D_1^2=16(1+i)^2y^4t^6=32iy^4t^6=4D_0.
\]
\end{proof}

The middle-branch reconstruction gives
\[
X_1=-\frac{D_1}{2}=-2(1+i)y^2t^3
\]
and
\begin{align}
X_2&=G_2+xX_1
=yt^2-2(1+i)xy^2t^3,\notag\\
X_3&=G_3+x^2X_1
=2xyt^2-t-2(1+i)x^2y^2t^3,\notag\\
X_4&=G_4+yX_1
=\left(-2-\frac{2i}{3}\right)y^3t^3.
\label{eq:X}
\end{align}
For
\[
\Delta=(1,x,x^2,y)^{\mathsf T}
\]
define the unpadded sweep
\[
S(\gamma,x,y,t)=X(x,y,t)+\gamma\Delta(x,y).
\]

\begin{proposition}\label{prop:sweep}
The sweep satisfies
\[
\det JS=\gamma^2.
\]
Moreover $X$ has generic rank three; for example
\[
\det\frac{\partial(X_1,X_2,X_3)}{\partial(x,y,t)}
=(-4+12i)y^3t^6.
\]
Thus $\Crit(S)=\{\gamma=0\}$ and the critical-value closure has dimension three.
\end{proposition}

\begin{proof}
Substitution of \eqref{eq:X} into the $4\times4$ Jacobian gives the polynomial identity $\det JS=\gamma^2$. The displayed $3\times3$ minor is nonzero, hence $\rank JX=3$ on a dense open subset. Since $\det JS=\gamma^2$, the set-theoretic critical locus is exactly $\{\gamma=0\}$; scheme-theoretically the determinant vanishes there with multiplicity two. Its image is $X(\A^3)$, whose closure has dimension three by generic rank three.
\end{proof}

\section{Elimination and the critical-value threefold}

Let $(z_1,z_2,z_3,z_4)$ be target coordinates and set
\[
R=z_1z_3-z_2^2.
\]
Write
\[
\kappa=2(1+i),
\qquad
\rho=1+2i,
\qquad
 d=-2-\frac{2i}{3}.
\]
From \eqref{eq:X},
\[
z_1=-\kappa y^2t^3,
\qquad
z_4=d y^3t^3.
\]
The $x$-terms cancel in $R$:
\begin{align*}
R
&=X_1X_3-X_2^2\\
&=\rho y^2t^4.
\end{align*}
Therefore
\begin{align*}
R^3z_4^2
&=\rho^3d^2y^{12}t^{18},\\
z_1^6
&=\kappa^6y^{12}t^{18}.
\end{align*}
Hence
\begin{equation}
R^3z_4^2=\lambda z_1^6,
\qquad
\lambda=\frac{\rho^3d^2}{\kappa^6}
=\frac{41-38i}{576}\neq0.
\label{eq:implicit}
\end{equation}
This is the promised elimination certificate; it requires no division by $y$ or $t$ and therefore holds on the entire parameter space.

\begin{proposition}\label{prop:closure}
The critical-value closure is exactly the irreducible threefold
\[
V=V(F)\subset\A^4,
\qquad
F=(z_1z_3-z_2^2)^3z_4^2-\lambda z_1^6.
\]
After a nonzero linear rescaling of $z_4$, one may equivalently take $\lambda=1$.
\end{proposition}

\begin{proof}
Equation \eqref{eq:implicit} shows that the image closure is contained in $V(F)$. It remains to prove irreducibility.

The polynomial
\[
R=z_1z_3-z_2^2
\]
is irreducible in $\C[z_1,z_2,z_3]$: it is primitive and linear in $z_3$, with relatively prime coefficient $z_1$ and constant term $-z_2^2$. Moreover $\gcd(R,z_1)=1$, so the coefficients $R^3$ and $-z_1^6$ of $F=R^3z_4^2-z_1^6$ have no common nonunit factor. Thus $F$ is primitive as a polynomial in $z_4$ over the UFD $\C[z_1,z_2,z_3]$. Over its fraction field, reducibility would require
\[
\frac{\lambda z_1^6}{R^3}
\]
to be a square. Its valuation along the irreducible prime $R$ is $-3$, which is odd; hence it is not a square. Gauss' lemma gives irreducibility of $F$.

By Proposition \ref{prop:sweep}, the image closure is irreducible of dimension three. The irreducible hypersurface $V(F)$ also has dimension three. An irreducible closed subset of an irreducible variety having the same dimension is the whole variety, so the critical-value closure equals $V(F)$.
\end{proof}

\section{The critical-value threefold is not a cylinder}

Because nonzero target rescaling does not affect the argument, normalize $\lambda=1$ and write
\[
F=R^3z_4^2-z_1^6,
\qquad
R=z_1z_3-z_2^2.
\]

\begin{lemma}\label{lem:sing}
The singular locus of $V$ is
\[
\Sing(V)=P_1\cup P_2,
\]
where
\[
P_1=\{z_1=z_2=0\},
\qquad
P_2=\{z_1=z_4=0\}.
\]
Their intersection is the line
\[
L=P_1\cap P_2=\{z_1=z_2=z_4=0\}.
\]
\end{lemma}

\begin{proof}
The partial derivatives are
\begin{align*}
F_{z_4}&=2R^3z_4,\\
F_{z_3}&=3R^2z_1z_4^2,\\
F_{z_2}&=-6R^2z_2z_4^2,\\
F_{z_1}&=3R^2z_3z_4^2-6z_1^5.
\end{align*}
At a singular point $F_{z_4}=0$, hence either $R=0$ or $z_4=0$.

If $R=0$, the equation $F=0$ gives $z_1=0$, and then $R=-z_2^2=0$, so $z_2=0$; this is $P_1$. Conversely every point of $P_1$ annihilates $F$ and all first derivatives.

If $z_4=0$, the equation $F=0$ gives $z_1=0$, yielding $P_2$, and again every point of $P_2$ is singular. Thus the singular locus is exactly $P_1\cup P_2$.
\end{proof}

\begin{lemma}\label{lem:mult}
For
\[
p_a=(0,0,a,0)\in L,
\]
one has
\[
\mult_{p_a}V=5\quad(a\neq0),
\qquad
\mult_0V=6.
\]
Thus the origin is the unique multiplicity-six point of $L$.
\end{lemma}

\begin{proof}
At $p_a$ with $a\neq0$, use local coordinates with $z_3=a+u_3$. Then
\[
R=az_1+z_1u_3-z_2^2
\]
has order one, with initial term $az_1$. Hence $R^3z_4^2$ has order five and nonzero degree-five initial term $a^3z_1^3z_4^2$, whereas $z_1^6$ has order six. Therefore the local multiplicity is five.

At the origin, $R$ has order two, so $R^3z_4^2$ has order eight, while $z_1^6$ has order six. Therefore $\mult_0V=6$.
\end{proof}

\begin{proposition}\label{prop:notcyl}
The threefold $V$ is not ambient-polynomially equivalent to a cylinder $Y\times\A^1\subset\A^3\times\A^1$.
\end{proposition}

\begin{proof}
Assume that a polynomial automorphism $\Psi$ of $\A^4$ sends $V$ onto a cylinder $Y\times\A^1$. Translation in the $\A^1$ coordinate defines an algebraic $\mathbb G_a$-action on $\A^4$ preserving the cylinder and having no global fixed point. Conjugating by $\Psi$ gives an algebraic $\mathbb G_a$-action on $\A^4$ preserving $V$ and likewise having no global fixed point.

Every ambient polynomial automorphism $\Theta$ preserving $V$ satisfies $\Theta(\Sing(V))=\Sing(V)$, because regularity of the local ring is preserved under isomorphism. The planes $P_1$ and $P_2$ are the two irreducible components of $\Sing(V)$. An algebraic action of a connected algebraic group on a finite discrete set is trivial: each orbit is the connected image of the group and hence consists of one point. Since $\mathbb G_a$ is connected and the set of irreducible components $\{P_1,P_2\}$ is finite, each component is therefore preserved setwise. Therefore their intersection $L$ is preserved setwise.

Local multiplicity is invariant under ambient polynomial automorphisms. By Lemma \ref{lem:mult}, the origin is the unique multiplicity-six point of $L$. Hence every element of the conjugated $\mathbb G_a$-action fixes the origin. This gives a global fixed point, contradicting the fixed-point-free translation action from which it was conjugated. Therefore no such ambient polynomial equivalence to a cylinder exists.
\end{proof}

\section{Excluding all triangular one-coordinate suspensions}

The determinant restrictions needed below are consequences, not assumptions. For a four-dimensional tangent sweep, Gao's Lemma 4.1 gives for the padded map
\[
F_0(x,\gamma_T,a,b)=\bigl(\gamma_T x,\,T(\gamma_T,a,b)\bigr)
\]
the identity
\[
\det J(F_0)=\gamma_T^2L+\gamma_T^3M.
\]
Here the exact padded/unpadded relation is immediate from the Jacobian: $T$ is independent of $x$, so expansion along the $x$-column gives
\[
\det J(F_0)=\gamma_T\,\det JT.
\]
Consequently the corresponding unpadded sweep $T:\A^3\to\A^3$ satisfies
\begin{equation}
\det JT=\gamma_T L+\gamma_T^2M
=\gamma_T(L+\gamma_T M),
\label{eq:gao4det}
\end{equation}
where $L,M$ are polynomials in the two non-$\gamma_T$ base parameters.

\begin{lemma}[Polynomial equivalence forces normalization]\label{lem:normalize}
Let $S$ be the sweep of Proposition \ref{prop:sweep}, with source coordinate $\gamma_S$ and
\[
\det JS=\gamma_S^2.
\]
Let
\[
\Sigma(u,v)=\bigl(T(u),Q(u,v)\bigr),
\qquad u=(\gamma_T,a,b),
\]
be any triangular one-coordinate suspension of a four-dimensional Gao sweep. If $S$ and $\Sigma$ are polynomially equivalent, then
\[
\det J\Sigma=B\gamma_T^2
\]
for some $B\in\C^\times$.
\end{lemma}

\begin{proof}
By symmetry of polynomial equivalence, write
\[
S=\Psi\circ\Sigma\circ\Phi
\]
for polynomial automorphisms $\Phi,\Psi$ of $\A^4$. Polynomial automorphisms have nonzero constant Jacobian determinant, so the chain rule and $\det JS=\gamma_S^2$ imply
\[
\det J\Sigma=Cp^2,
\qquad
p:=\gamma_S\circ\Phi^{-1},
\qquad C\in\C^\times.
\]
The polynomial $p$ is a coordinate polynomial, hence irreducible.  Indeed the automorphism $\Phi^{-1}$ identifies $\C[p,p_2,p_3,p_4]$ with the full polynomial ring for suitable companion coordinates; equivalently, the quotient by $(p)$ is isomorphic to a polynomial ring in three variables and is therefore a domain.  Thus $(p)$ is prime, in particular $p$ is irreducible.

Because $\Sigma=(T,Q)$ is triangular and $T$ is independent of $v$,
\[
\det J\Sigma=(\partial_vQ)\det JT.
\]
By \eqref{eq:gao4det}, $\gamma_T$ divides $\det JT$, hence
\[
\gamma_T\mid\det J\Sigma=Cp^2.
\]
The ring $\C[\gamma_T,a,b,v]$ is a UFD and the coordinate $\gamma_T$ is irreducible, hence prime. Therefore $\gamma_T\mid p^2$ implies $\gamma_T\mid p$. Since $p$ and $\gamma_T$ are both irreducible, they are associates:
\[
p=A\gamma_T,
\qquad A\in\C^\times.
\]
Consequently
\[
\det J\Sigma=CA^2\gamma_T^2,
\]
which is the asserted normalization.
\end{proof}

\begin{lemma}[Triangular suspension exhaustiveness]\label{lem:det}
Let
\[
\Sigma(u,v)=\bigl(T(u),Q(u,v)\bigr),
\qquad u=(\gamma_T,a,b),
\]
be a triangular one-coordinate suspension satisfying
\[
\det J\Sigma=B\gamma_T^2,
\qquad B\in\C^\times.
\]
Write the general four-dimensional Gao sweep determinant as in \eqref{eq:gao4det}. Then exactly one of the following two cases holds for some $A\in\C^\times$:
\begin{enumerate}[label=(\Alph*)]
\item $M=0$, $L=A$, hence
\[
\det JT=A\gamma_T,
\qquad
Q=\frac BA\,\gamma_T v+h(u);
\]
\item $L=0$, $M=A$, hence
\[
\det JT=A\gamma_T^2,
\qquad
Q=\frac BA\,v+h(u).
\]
\end{enumerate}
In particular these two determinant modes are exhaustive; they are not part of the definition.
\end{lemma}

\begin{proof}
Because $T$ is independent of $v$, the Jacobian matrix is block triangular:
\[
J\Sigma=
\begin{pmatrix}
JT&0\\
\partial_uQ&\partial_vQ
\end{pmatrix},
\]
so
\[
\det J\Sigma=(\partial_vQ)\det JT.
\]
Using \eqref{eq:gao4det} and cancelling the nonzero factor $\gamma_T$ in the integral domain $\C[\gamma_T,a,b,v]$ gives
\begin{equation}
(\partial_vQ)(L+\gamma_T M)=B\gamma_T.
\label{eq:ufd}
\end{equation}
The polynomial ring is a UFD and $\gamma_T$ is irreducible. Hence the nonzero factor $L+\gamma_T M$ divides $B\gamma_T$. Since $B$ is a unit, every nonunit irreducible factor of $L+\gamma_T M$ must be associated to $\gamma_T$. Moreover $L+\gamma_T M$ has degree at most one in $\gamma_T$. Therefore it is either a unit or an associate of $\gamma_T$.

If it is a unit, units in $\C[\gamma_T,a,b,v]$ are nonzero constants, so $M=0$ and $L=A\in\C^\times$. Equation \eqref{eq:ufd} gives $\partial_vQ=(B/A)\gamma_T$, and integration in $v$ yields
\[
Q=(B/A)\gamma_T v+h(u).
\]
If $L+\gamma_T M$ is associated to $\gamma_T$, then $L=0$ and $M=A\in\C^\times$. Equation \eqref{eq:ufd} gives $\partial_vQ=B/A$, hence
\[
Q=(B/A)v+h(u).
\]
The two alternatives are mutually exclusive and exhaustive.
\end{proof}

\begin{lemma}[Critical-value invariant]\label{lem:equiv}
If $S_2=\Psi\circ S_1\circ\Phi$ with polynomial automorphisms $\Phi,\Psi$ of $\A^4$, then
\[
\overline{\CV(S_2)}=\Psi\bigl(\overline{\CV(S_1)}\bigr).
\]
In particular the dimension of the critical-value closure and the property of being ambient-polynomially equivalent to a cylinder are invariants of source-target polynomial equivalence.
\end{lemma}

\begin{proof}
Polynomial automorphisms have everywhere invertible Jacobian matrices. By the chain rule,
\[
\Crit(S_2)=\Phi^{-1}(\Crit(S_1)).
\]
Applying $S_2$ gives
\[
\CV(S_2)=\Psi(\CV(S_1)),
\]
and taking Zariski closures yields the formula. Dimension is preserved by $\Psi$, and cylinder-equivalence is preserved by composition with ambient polynomial automorphisms.
\end{proof}

\begin{proposition}\label{prop:exclude}
The sweep $S$ of Proposition \ref{prop:sweep} is not polynomially equivalent to any triangular one-coordinate suspension of a four-dimensional Gao sweep.
\end{proposition}

\begin{proof}
Suppose for contradiction that $S$ is polynomially equivalent to a triangular one-coordinate suspension $\Sigma=(T,Q)$. Lemma \ref{lem:normalize} forces
\[
\det J\Sigma=B\gamma_T^2,
\qquad B\in\C^\times.
\]
Lemma \ref{lem:det} then forces its determinant structure into exactly one of cases (A) or (B).

In case (A),
\[
Q=\alpha\gamma_T v+h(u),\qquad \alpha\neq0.
\]
The set-theoretic critical locus is $\gamma_T=0$. On that locus the $v$-dependence disappears from $Q$, so the critical-value set is parameterized by the two remaining base parameters. Hence
\[
\dim\overline{\CV(\Sigma)}\le2.
\]
Our sweep has a three-dimensional critical-value closure by Proposition \ref{prop:sweep}. Lemma \ref{lem:equiv} excludes polynomial equivalence.

In case (B),
\[
Q=\alpha v+h(u),\qquad \alpha\neq0.
\]
On the set-theoretic critical locus $\gamma_T=0$, the coordinate $v$ is free and the last target coordinate ranges over all of $\A^1$ for every base critical value. Thus
\[
\overline{\CV(\Sigma)}
=
\overline{\CV(T)}\times\A^1,
\]
a cylinder. Proposition \ref{prop:notcyl} and Lemma \ref{lem:equiv} again exclude polynomial equivalence.

Lemmas \ref{lem:normalize} and \ref{lem:det} prove that any triangular one-coordinate suspension polynomially equivalent to $S$ would be normalized automatically and would fall into exactly one of these two cases. Both cases are impossible. Hence no triangular one-coordinate suspension of a four-dimensional Gao sweep is polynomially equivalent to $S$.
\end{proof}

\begin{proof}[Proof of Theorem \ref{thm:main}]
Proposition \ref{prop:master} gives an exact polynomial middle-branch solution. Proposition \ref{prop:exclude} proves that its associated sweep is not polynomially equivalent to any triangular one-coordinate suspension of a four-dimensional Gao sweep. This is precisely the assertion.
\end{proof}

\begin{proof}[Proof of Corollary \ref{cor:gao}]
Under the stated interpretation of Gao's phrase ``equivalent to a suspension of a four-dimensional sweep,'' Theorem \ref{thm:main} supplies the polynomial solution whose existence is asked for in Open Problem 4.8. Hence the conjectured rigidity fails.
\end{proof}

\section{Scope and verification}

The main theorem is intentionally precise about the word ``suspension.''  As emphasized in Remark \ref{rem:scope}, it excludes the formally defined triangular one-coordinate Gao-suspension class and does not assert a classification of every possible informal use of that word. Definition 2 imposes only the natural triangular one-coordinate form with a four-dimensional Gao base sweep; it imposes no determinant normalization. Lemma \ref{lem:normalize} proves that polynomial equivalence to our square-Jacobian sweep would itself force $\det J\Sigma\sim\gamma_T^2$. Lemma \ref{lem:det}, using Gao's general four-dimensional sweep formula (Lemma 4.1), then proves that the two possibilities $\det JT\sim\gamma_T$ and $\det JT\sim\gamma_T^2$ are exhaustive. Gao's Theorem 4.6 supplies a concrete instance of case (A). Remark 4.7 motivates the graph-coordinate type of collapse but is not used as a formal classification theorem for case (B).

All witness identities in Propositions \ref{prop:master} and \ref{prop:sweep}, the rank-three minor, and the implicit relation \eqref{eq:implicit} have been checked in exact symbolic arithmetic. The accompanying verification script is intended as a reproducibility aid; the proof above does not depend on randomized computation.

\begin{thebibliography}{9}
\bibitem{Gao2026}
S. Gao,
\emph{Counterexamples to the Jacobian conjecture in dimensions greater than two},
arXiv:2608.00222v1 [math.AG], 31 July 2026.
Relevant items: Lemma 4.1, equation (12), Theorem 4.6, Remark 4.7, Open Problem 4.8.
\end{thebibliography}

\end{document}
